% from main.tex
%%%%%%%%%%  scan idx 248  |  folio 235  %%%%%%%%%%
\emph{Preuve de a).} D'après (1.1 (iii)), on doit montrer ceci : supposons \(f_*(E)\) \(a\)-pseudo-cohérent et acyclique en degré \(>b\), et \(g_*(F)\) \(c\)-pseudo-cohérent et acyclique en degré \(>d\), alors \(h_*(E\otimes_Y^{\mathbf L}F)\) est \(m\)-pseudo-cohérent et acyclique en degré \(>n\). Or on a un isomorphisme canonique
\[
  f_*(E)\otimes^{\mathbf L}F\xrightarrow{\sim} f_*(E\otimes_Y^{\mathbf L}F)
\]
(cas particulier de l'isomorphisme (*) de la preuve de (1.8)). Comme le foncteur \(g_*(.\otimes^{\mathbf L}F)\) transforme faisceaux localement libres de type fini en complexes \(c\)-pseudo-cohérents, et complexes acycliques en degré \(>0\) en complexes acycliques en degré \(>d\), l'assertion résulte donc de (I 2.13).

\emph{Preuve de b).} On doit prouver l'assertion suivante : supposons \(f_*(E)\) \(a\)-pseudo-cohérent et acyclique en degré \(>b\), et \(g_*(\mathcal O_Y)\) \((a+1-b)\)-pseudo-cohérent, alors \(f_*(E)\) est \(a\)-pseudo-cohérent. L'assertion portant sur \(f_*(E)\), on peut supposer \(X=Y\) et \(f=\Id_X\). On procède alors comme dans la démonstration de (1.1.1), en prouvant par récurrence descendante que \(E\) est pseudo-cohérent en degré \(\ge a\). On suppose déjà prouvé que \(E\) est \((a+1)\)-pseudo-cohérent, d'où (I 2.10), quitte à se localiser, un triangle distingué
\[
E'\longrightarrow E\longrightarrow E''\longrightarrow E'[1]
\]
où \(E'^i\) est localement libre pour \(i\in[a+1,b]\), \(E'^i=0\) pour \(i\notin[a+1,b]\), et \(E''^i=0\) pour \(i>a\). L'hypothèse que \(g_*(\mathcal O_Y)\) est \((a+1-b)\)-pseudo-cohérent implique alors (I 2.11 a)) que \(g_*(E')\) est \((a+1)\)-pseudo-cohérent. Comme \(g_*(E)\) est \(a\)-pseudo-cohérent, il résulte de (I 2.6) que \(g_*(E'')\) est \(a\)-pseudo-cohérent. On conclut grâce à (I 2.10 b)) et au fait qu'un \(\mathcal O_Y\)-Module \(M\) est de type fini si et seulement si \(g_*(M)\) l'est.

On déduit trivialement de (1.11) les corollaires que voici :

